\section{Meta-Policy for Test Initiation}
\label{sec:meta_policy}

Not every restriction should be tested.  Controlled relaxation
carries costs: the damage tolerance $\damagebound$ is a real (if
bounded) risk, the test consumes monitoring resources, and
concurrent tests raise the chance of correlated failures.  The
meta-policy determines when a test is worth running.

\subsection{When to test}

\begin{definition}[Test Initiation Criteria]
\label{def:test_initiation}
A controlled relaxation test for risk-restriction pair
$(R, \Xirestrict)$ is initiated when all of the following hold.
\begin{itemize}
\item[(i)] \textbf{Stale restriction.}  $\Xirestrict$ has
  persisted for at least $T_{\mathrm{stale}}$ time steps without
  new supporting evidence for $R$.  ``New supporting evidence''
  means any observation that increased $\pR$ via a Bayesian
  update (not merely the passage of time under restriction, which
  provides no evidence by
  Proposition~\ref{prop:asym_evidence}).
\item[(ii)] \textbf{Favourable cost-benefit.}  The expected
  information value of the test exceeds the expected damage cost,
  \begin{equation}
  \label{eq:iv_criterion}
  \IV(T)
     = \mathbb{E}_{O}\bigl[
         \KL(\pR^{\mathrm{post}} \,\|\, \pR^{\mathrm{prior}})
        \bigr]
     > \damagebound \cdot \pR(t),
  \end{equation}
  where the expectation is over the test outcome weighted by its
  probability and $\damagebound \cdot \pR(t)$ is the expected
  damage (tolerance times probability the risk materializes).
\item[(iii)] \textbf{Scope isolation.}  No other controlled
  relaxation test is currently running on a scope $\scope'$ that
  overlaps with $\scope$.
\item[(iv)] \textbf{User approval (when applicable).}  If the
  damage tolerance $\damagebound$ exceeds a user-configured
  threshold $\damagebound_{\mathrm{user}}$, the test requires
  explicit user approval before initiation.  This is the safety
  valve: the framework can initiate low-damage tests
  autonomously but must ask the user before high-damage tests.
  The threshold $\damagebound_{\mathrm{user}}$ is analogous to
  the self-trust pause threshold in the GFM harness
  constitution.
\item[(v)] \textbf{Empirical exercise consistency.}  The
  empirically observed exercise frequency $\hat{p}_{\mathrm{ex}}$
  over the running test sequence (the fraction of test windows
  in which detection-side activation indicators register, used
  as a proxy for actual capability exercise) must remain within a
  configured tolerance of the modelled $p_{\mathrm{ex}}$
  (Assumption~\ref{ass:T5}).  Test sequences for which
  $\hat{p}_{\mathrm{ex}}$ falls below modelled $p_{\mathrm{ex}}$
  by more than the tolerance are treated as test-design
  failures: the test outcomes are discarded, the restriction
  returns to its prior state, and the meta-policy logs the
  failure for review.  This guards against the upward
  $p_{\mathrm{ex}}$-misspecification failure mode of
  Remark~\ref{rem:pex_robustness}.
\end{itemize}
\end{definition}

\begin{remark}[Connection to structural discovery value]
\label{rem:structural_preference}
Criterion~(ii) is a specialization of the structural discovery
value of~\citep[Definition~5,
Proposition~2]{lasser2026horizon}: a structural discovery that
would resolve the risk claim has information value
$\IV_{\mathrm{struct}}$; a controlled relaxation test has
information value $\IV(T)$.  The framework should prefer
whichever is cheaper: if a structural argument can resolve the
claim without exercise (for example, proving the pathway is
mechanically infeasible), that is strictly better than any
controlled relaxation test.  Controlled relaxation is the
fallback when no structural resolution is available---the
empirical arbitration mechanism for claims that resist formal
resolution.
\end{remark}

\subsection{Test parameter optimization}

\begin{proposition}[Optimal Test Duration]
\label{prop:tau_opt}
For fixed scope $\scope$ and damage tolerance $\damagebound$,
the test duration $\testdur$ that maximizes the expected
information value per unit of risk exposure is
\begin{equation}
\label{eq:tau_opt}
\testopt = \arg\max_{\tau}
  \frac{\IV(T(\tau))}{\damagebound \cdot \Pr(\text{contraction in } \tau \text{ steps})}.
\end{equation}
Under the exponential detection model
$\beta(\tau) = (1 - \deltaagg)^{\tau}$
(Assumption~\ref{ass:D2}(a)) and persistent activation
(Assumption~\ref{ass:D0}),
\begin{equation}
\label{eq:tau_opt_explicit}
\testopt \approx \left\lceil \frac{-1}{\log(1 - \deltaagg)} \right\rceil,
\end{equation}
the time constant of the exponential detection process---the
number of steps for the miss probability to decay by a factor
of $e$.  For $\deltaagg \ll 1$ this is approximately
$1/\deltaagg$; for $\deltaagg$ close to $1$, $\testopt = 1$
suffices.
\end{proposition}

\begin{proof}[Proof sketch]
The information value $\IV$ scales as $|\log \beta(\tau)|$ for
the no-contraction case, which grows linearly in $\tau$.  The
risk exposure grows as
$1 - (1 - \pR^{\mathrm{true}} \deltaagg)^{\tau}
 \approx \pR^{\mathrm{true}} \deltaagg \tau$
for small $\tau$.  The ratio $\IV/\mathrm{risk}$ peaks when the
marginal information gain equals the marginal risk increase,
which occurs at $\tau \approx 1/\deltaagg$.  Running longer
provides diminishing returns ($\beta$ is already small); running
shorter wastes the test setup cost.
\end{proof}

\subsection{Scheduling multiple tests}

\begin{definition}[Test Schedule]
\label{def:test_schedule}
A \emph{test schedule}
$\Sigma = \{(R_i, \Xirestrict_i, \scope_i, \tau_{\mathrm{test},i},
   \damagebound_i)\}_{i=1}^{N}$
is a set of controlled relaxation tests to be run.  The schedule
is \emph{feasible} if:
\begin{itemize}
\item[(i)] all scopes $\scope_i$ are pairwise non-overlapping;
\item[(ii)] the aggregate damage tolerance
  $\sum_i \damagebound_i \leq \damagebound_{\mathrm{total}}$ is
  within the framework's total damage budget;
\item[(iii)] tests are prioritized by their
  information-value-to-cost ratio
  $\IV(T_i) / \damagebound_i$ (highest ratio first).
\end{itemize}
\end{definition}

\begin{remark}[Scheduling is a constrained knapsack]
\label{rem:knapsack}
The scheduling problem is a variant of the knapsack problem
with isolation constraints (non-overlapping scopes).  For small
numbers of concurrent tests, greedy selection by
$\IV / \damagebound$ with scope-conflict checking is
sufficient; for large-scale deployment, integer-programming
formulations are available.
\end{remark}
